
\documentclass[11pt,a4paper]{article}
\usepackage[T1]{fontenc}
\usepackage{lmodern}
\usepackage{amsmath,amssymb,amsfonts,mathtools}
\usepackage{graphicx,booktabs,longtable,array,tabularx,adjustbox}
\usepackage{geometry,microtype,fancyhdr,lastpage,etoolbox,enumitem}
\usepackage{newunicodechar,textcomp}
\usepackage{hyperref}
\hypersetup{hidelinks,pdfcreator={XeLaTeX / QM1 academic edition}}
\geometry{left=20mm,right=20mm,top=22mm,bottom=23mm}
\setcounter{secnumdepth}{2}
\setlength{\parindent}{0pt}
\setlength{\parskip}{4pt plus 1pt}
\setlength{\emergencystretch}{2em}
\setlength{\tabcolsep}{3pt}
\renewcommand{\arraystretch}{1.16}
\setlist{nosep,topsep=3pt,itemsep=2pt,leftmargin=7mm}
\AtBeginEnvironment{longtable}{\footnotesize\setlength{\tabcolsep}{2.2pt}}
\newcommand{\tightlist}{\setlength{\itemsep}{0pt}\setlength{\parskip}{0pt}}
\providecommand{\Cov}{\operatorname{Cov}}
\providecommand{\Var}{\operatorname{Var}}
\providecommand{\Corr}{\operatorname{Corr}}
\providecommand{\MAD}{\operatorname{MAD}}
\providecommand{\E}{\mathbb{E}}
\providecommand{\R}{\mathbb{R}}
\newunicodechar{—}{\textemdash}
\newunicodechar{–}{\textendash}
\newunicodechar{€}{\texteuro}
\newunicodechar{≤}{\ensuremath{\leq}}
\newunicodechar{≥}{\ensuremath{\geq}}
\newunicodechar{×}{\ensuremath{\times}}
\newunicodechar{÷}{\ensuremath{\div}}
\newunicodechar{Σ}{\ensuremath{\Sigma}}
\newunicodechar{σ}{\ensuremath{\sigma}}
\newunicodechar{μ}{\ensuremath{\mu}}
\newunicodechar{Δ}{\ensuremath{\Delta}}
\newunicodechar{∑}{\ensuremath{\sum}}
\newunicodechar{≈}{\ensuremath{\approx}}
\newunicodechar{√}{\ensuremath{\surd}}
\newunicodechar{≠}{\ensuremath{\ne}}
\newunicodechar{→}{\ensuremath{\to}}
\newunicodechar{∞}{\ensuremath{\infty}}
\newunicodechar{∈}{\ensuremath{\in}}
\newunicodechar{−}{\ensuremath{-}}
\newunicodechar{²}{\ensuremath{^{2}}}
\newunicodechar{³}{\ensuremath{^{3}}}
\newunicodechar{₀}{\ensuremath{_{0}}}
\newunicodechar{₁}{\ensuremath{_{1}}}
\newunicodechar{₂}{\ensuremath{_{2}}}
\newunicodechar{₃}{\ensuremath{_{3}}}
\newunicodechar{₄}{\ensuremath{_{4}}}
\newunicodechar{ₜ}{\ensuremath{_{t}}}
\newunicodechar{ₚ}{\ensuremath{_{p}}}
\newunicodechar{ₙ}{\ensuremath{_{n}}}
\newunicodechar{ₓ}{\ensuremath{_{x}}}
\newunicodechar{ᵢ}{\ensuremath{_{i}}}
\newunicodechar{ŷ}{\ensuremath{\hat y}}
\newunicodechar{Ŷ}{\ensuremath{\widehat Y}}
\newunicodechar{ₘ}{\ensuremath{_{m}}}
\newunicodechar{ₑ}{\ensuremath{_{e}}}
\newunicodechar{ₐ}{\ensuremath{_{a}}}
\newunicodechar{ₖ}{\ensuremath{_{k}}}
\newunicodechar{ₛ}{\ensuremath{_{s}}}
\newunicodechar{ₗ}{\ensuremath{_{l}}}
\newunicodechar{ₕ}{\ensuremath{_{h}}}
\newunicodechar{ₒ}{\ensuremath{_{o}}}
\newunicodechar{₊}{\ensuremath{_{+}}}
\newunicodechar{₋}{\ensuremath{_{-}}}
\newunicodechar{₌}{\ensuremath{_{=}}}
\newunicodechar{₍}{\ensuremath{_{(}}}
\newunicodechar{₎}{\ensuremath{_{)}}}
\newunicodechar{⁻}{\ensuremath{^{-}}}
\newunicodechar{⁺}{\ensuremath{^{+}}}
\newunicodechar{⁰}{\ensuremath{^{0}}}
\newunicodechar{¹}{\ensuremath{^{1}}}
\newunicodechar{⁴}{\ensuremath{^{4}}}
\newunicodechar{⁵}{\ensuremath{^{5}}}
\newunicodechar{⁶}{\ensuremath{^{6}}}
\newunicodechar{⁷}{\ensuremath{^{7}}}
\newunicodechar{⁸}{\ensuremath{^{8}}}
\newunicodechar{⁹}{\ensuremath{^{9}}}
\makeatletter
\renewcommand{\section}{\@startsection{section}{1}{0pt}{2.0ex plus .7ex minus .3ex}{1.0ex}{\large\bfseries}}
\renewcommand{\subsection}{\@startsection{subsection}{2}{0pt}{1.8ex plus .5ex minus .2ex}{.7ex}{\normalsize\bfseries}}
\renewcommand{\subsubsection}{\@startsection{subsubsection}{3}{0pt}{1.3ex plus .4ex minus .2ex}{.5ex}{\normalsize\bfseries}}
\makeatother
\pagestyle{fancy}
\fancyhf{}
\renewcommand{\headrulewidth}{0.35pt}
\renewcommand{\footrulewidth}{0pt}
\fancyhead[L]{\small\textsc{KEDGE Business School --- QM1}}
\fancyhead[R]{\small\textsc{2026--2027}}
\fancyfoot[L]{\small Badr ABOUFARASSE, PhD}
\fancyfoot[R]{\small\thepage\ /\ \pageref{LastPage}}
\setlength{\headheight}{14pt}
\setkeys{Gin}{keepaspectratio}
\begin{document}

\begin{center}
{\small\textsc{KEDGE Business School \quad EBP1 \quad Quantitative Methods 1}}\\[6pt]
{\Large\bfseries Worksheet 1 — Presenting and organising data}\\[6pt]
{\normalsize\bfseries Exercise statements}\\[6pt]
{\small Badr ABOUFARASSE, PhD \quad --- \quad 2026--2027}
\end{center}
\vspace{2mm}\hrule\vspace{6mm}

\hypertarget{core-applications}{%
\section{Core applications}\label{core-applications}}

Master the essential definitions and calculations from the course.

\emph{Core application --- 10 min}

\hypertarget{exercise-1-population-sample-and-variable-types}{%
\subsection{Exercise 1 -- Population, sample and variable types}\label{exercise-1-population-sample-and-variable-types}}

A store wants to study all customers who made a purchase in September. It surveys a sample of 80 customers and records:

\begin{longtable}[]{@{}p{0.32\linewidth}p{0.62\linewidth}@{}}
\toprule\noalign{}
Variable & Example \\
\midrule\noalign{}
\endhead
\bottomrule\noalign{}
\endlastfoot
Payment method & card, cash, mobile \\
Number of items purchased & 1, 2, 3, \ldots{} \\
Time spent in store & minutes \\
Satisfaction level & very dissatisfied to very satisfied \\
\end{longtable}

\begin{enumerate}
\item
  Identify the target population, the sample and the statistical unit.
\item
  Classify each variable as nominal qualitative, ordinal qualitative, discrete quantitative or continuous quantitative.
\item
  Choose an appropriate graph for each variable.
\item
  For which variables do cumulative frequencies make sense? Explain.
\item
  Explain why a large sample is not necessarily representative.
\end{enumerate}

\emph{Core application --- 20 min}

\hypertarget{exercise-2-discrete-distribution}{%
\subsection{Exercise 2 -- Discrete distribution}\label{exercise-2-discrete-distribution}}

\begin{longtable}[]{@{}lllllll@{}}
\toprule\noalign{}
{\(x_i\)} & 0 & 1 & 2 & 3 & 4 & 5 \\
\midrule\noalign{}
\endhead
\bottomrule\noalign{}
\endlastfoot
{\(n_i\)} & 3 & 7 & 12 & 10 & 6 & 2 \\
\end{longtable}

\begin{enumerate}
\item
  Check the total count {\(N\)}.
\item
  Calculate relative frequencies {\(f_i=n_i/N\)} and increasing cumulative frequencies.
\item
  Draw a stick/bar diagram.
\item
  Construct the cumulative distribution function {\(F(x)=P(X\leq x)\)}.
\item
  Find the proportion of customers who placed at most three orders.
\end{enumerate}

\emph{Core application --- 15 min}

\hypertarget{extra-1-cross-tabulation-and-conditional-percentages}{%
\subsection{Extra 1 --- Cross-tabulation and conditional percentages}\label{extra-1-cross-tabulation-and-conditional-percentages}}

Among 120 students, 50 use public transport, including 30 women. Of the 70 other students, 35 are women.

\begin{enumerate}
\tightlist
\item
  Build the gender × transport table.
\item
  Compute the public-transport rate within each gender.
\item
  State whether an association appears present.
\end{enumerate}

\textbf{Method hint:}

Conditional percentages use the total of the conditioning group as denominator.

\hypertarget{assessment-and-exam-practice}{%
\section{Assessment and exam practice}\label{assessment-and-exam-practice}}

Practise independently with assessment and final-exam style questions.

\emph{Assessment / exam practice --- 40 min}

\hypertarget{exercise-3-continuous-variable-with-unequal-class-widths}{%
\subsection{Exercise 3 -- Continuous variable with unequal class widths}\label{exercise-3-continuous-variable-with-unequal-class-widths}}

\begin{longtable}[]{@{}llllllllll@{}}
\toprule\noalign{}
Class & {\([10,14[\)} & {\([14,18[\)} & {\([18,26[\)} & {\([26,30[\)} & {\([30,35[\)} & {\([35,40[\)} & {\([40,45[\)} & {\([45,50[\)} & {\([50,55[\)} \\
\midrule\noalign{}
\endhead
\bottomrule\noalign{}
\endlastfoot
{\(n_i\)} & 4 & 12 & 32 & 32 & 60 & 35 & 25 & 5 & 5 \\
\end{longtable}

\begin{enumerate}
\item
  For each class, calculate width {\(a_i\)}, relative frequency {\(f_i\)} and count density {\(d_i=n_i/a_i\)}.
\item
  Draw the correct histogram and explain why raw counts cannot be used as heights when class widths differ.
\item
  Calculate increasing and decreasing cumulative counts and the corresponding cumulative relative frequencies.
\item
  Draw the cumulative curves.
\item
  Find the number and proportion willing to pay less than 35 euros, then at least 35 euros.
\end{enumerate}

\emph{Assessment / exam practice --- 35 min}

\hypertarget{extra-3-continuous-assessment-case-willingness-to-pay}{%
\subsection{Extra 3 --- Continuous-assessment case: willingness to pay}\label{extra-3-continuous-assessment-case-willingness-to-pay}}

\textbf{Case study --- 12 marks.} A food retailer asks 200 customers for the maximum price they would pay for a sustainable lunch box.

\begin{longtable}[]{@{}lllllll@{}}
\toprule\noalign{}
Maximum price (€) & {[}20,25{[} & {[}25,30{[} & {[}30,40{[} & {[}40,50{[} & {[}50,60{[} & {[}60,80{[} \\
\midrule\noalign{}
\endhead
\bottomrule\noalign{}
\endlastfoot
Customers & 10 & 20 & 50 & 60 & 40 & 20 \\
\end{longtable}

\begin{enumerate}
\tightlist
\item
  Identify the population, statistical unit, variable and variable type. \emph{(2 marks)}
\item
  Calculate class widths, relative frequencies and frequency densities. \emph{(3 marks)}
\item
  Explain how to construct the correct histogram and identify the modal class. \emph{(2 marks)}
\item
  Calculate ascending and descending cumulative counts. Give the coordinates needed to draw both cumulative curves. \emph{(3 marks)}
\item
  Find the highest price that at least 60\% of surveyed customers would be willing to pay. Explain the percentile used. \emph{(2 marks)}
\end{enumerate}

\textbf{Method hint:}

Unequal-width classes require frequency density. For the pricing decision, translate ``at least 60\% can pay'' into an ascending cumulative percentage.

\emph{Assessment / exam practice --- 30 min}

\hypertarget{extra-4-exam-style-synthesis-reporting-customer-satisfaction}{%
\subsection{Extra 4 --- Exam-style synthesis: reporting customer satisfaction}\label{extra-4-exam-style-synthesis-reporting-customer-satisfaction}}

\textbf{Independent case --- 10 marks.} Two stores collect satisfaction ratings on an ordered scale from 1 (very dissatisfied) to 5 (very satisfied).

\begin{longtable}[]{@{}lllllll@{}}
\toprule\noalign{}
Rating & 1 & 2 & 3 & 4 & 5 & Total \\
\midrule\noalign{}
\endhead
\bottomrule\noalign{}
\endlastfoot
Store A & 6 & 10 & 18 & 36 & 30 & 100 \\
Store B & 4 & 12 & 30 & 34 & 20 & 100 \\
\end{longtable}

\begin{enumerate}
\tightlist
\item
  Identify the variable and its measurement scale. \emph{(1 mark)}
\item
  Choose and justify one graph for comparing the two stores. Describe precisely how it should be constructed. \emph{(2 marks)}
\item
  Compute, for each store, the percentages rated 4 or 5 and the cumulative percentage rated at most 2. \emph{(3 marks)}
\item
  A manager says: ``Store A is 50\% better because it has 30 ratings of 5 versus 20.'' Criticise this claim. \emph{(2 marks)}
\item
  Write a three-sentence statistical conclusion, including one limitation of the survey. \emph{(2 marks)}
\end{enumerate}

\textbf{Method hint:}

Keep the order of the categories, compare percentages because the store totals could differ in another dataset, and separate description from causation.

\hypertarget{optional-advanced-practice}{%
\section{Optional advanced practice}\label{optional-advanced-practice}}

Explore additional interpretations and professional applications.

\emph{Advanced practice --- 15 min}

\hypertarget{extra-2-data-quality-and-communication-protocol}{%
\subsection{Extra 2 --- Data-quality and communication protocol}\label{extra-2-data-quality-and-communication-protocol}}

A CRM extract contains duplicate customers, missing ages coded 999, inconsistent dates and revenue recorded in euros and thousands of euros. Propose a reproducible six-step protocol before producing a management dashboard.

\textbf{Method hint:}

Separate validation, cleaning, documentation and communication.

\end{document}
